% \iffalse meta-comment
%
% Copyright (C) 2017- by Yanshuo Chu <yanshuoc@gmail.com>
%
% This file may be distributed and/or modified under the
% conditions of the LaTeX Project Public License, either version 1.3a
% of this license or (at your option) any later version.
% The latest version of this license is in:
%
% http://www.latex-project.org/lppl.txt
%
% and version 1.3a or later is part of all distributions of LaTeX
% version 2004/10/01 or later.
%
% \fi
%
% \section{词条}
%
% \changes{v3.2a}{2026/09/13}{固定词汇改成 \option{term} 族的常量键，从
%   \file{hit-defaults.dtx} 按关切拆出来：原先一半走 \cs{ctexset}、一半靠
%   \cs{cs\_new:Npn} 直接定义命令，值里还嵌着 \cs{hit@if@option@TF}；现在按变体的
%   差别一律写成默认值加覆盖表，两条路都由 \file{hit-key-engine.dtx} 提供，配置只出取值}
% \changes{v3.2a}{2026/09/13}{取值照范例与新规范：校区名只留校区两个字；成果章一律叫
%   「攻读□士学位期间取得创新性成果」；本科原创性声明三个校区统一到本部这一份并按
%   新规范加 AI 工具那一句；终稿独有的元信息字段给范例占位形的回落值；报告页眉整句、
%   报告封面按变体分岔的取值、内封的学位说明、答辩决议那一页的表头文案（\issue[247]）
%   都成词条；深圳本科报告首页的词条中英各一套}
% \changes{v3.2a}{2026/08/29}{日期：排法迁到 \file{src/engine/hit-date.dtx}，默认值改成
%   东八区今天的年月，精度跟着值走（规范四处落点都是年月；时区原先跟编译机器走，国外
%   编译会差一天），深圳本科按 \issue{220} 单独给到日}
% \changes{v3.2a}{2026/08/27}{学生类别的两条文本成常量 \option{degree-type-academic}
%   与 \option{degree-type-professional}，名字照 iota；参考文献与索引的标题、英文图表名
%   收进 \option{term} 族}
%
% 固定词汇的取值，三段：两档共用一段，终稿一段，开题与中期一段。
% 段条件挂在各段段首，末尾撤销。
%
%    \begin{macrocode}
%<*hithesis-cfg>
%    \end{macrocode}
%
% 校区名。值只有校区两个字，「（威海）」的括号、「深圳校区」的「校区」都在用的地方
% 拼——同一个校区名在封面上带括号、在深圳本科页眉上跟着「校区」，那是版式的事，
% 不是这个值的事。三个类共用一份。
%    \begin{macrocode}
\ExplSyntaxOff
\hit@presetup[campus=harbin]{ term = { base-campus = {哈尔滨} } }
\hit@presetup[campus=weihai]{ term = { base-campus = {威海} } }
\hit@presetup[campus=shenzhen]{ term = { base-campus = {深圳} } }
\hit@presetup{ term = { base-campus-full = {<base-campus>校区} } }
\hit@presetup{ term = { cover-bottommark = {哈尔滨工业大学教务处制} } }
\ExplSyntaxOn
%    \end{macrocode}
%
%    \begin{macrocode}
\hit@presetup@scope{stage=final}
%    \end{macrocode}
%
%    \begin{macrocode}
\ExplSyntaxOff
\hit@presetup[category=hass]{ term = {
  abstract-in-english-toc-zh  = {Abstract(In Chinese)},
  abstract-in-english-toc-en  = {Abstract(In English)},
  acknowledgements-heading-en = {Acknowledgments},
  declarations-heading-en     = {Statement of copyright and letter of authorization},
} }
\ExplSyntaxOn
%    \end{macrocode}
%
%    \begin{macrocode}
\ExplSyntaxOff
%    \end{macrocode}
%
% 中英文的固定词汇。\option{chapter-name} 到 \option{indexname} 那几个归 \pkg{ctex} 管，
% 设值时转给 \cs{ctexset}；其余几个是对外可见的命令名，设值时直接写进同名命令。
% 两类都在 \file{hit-key-engine.dtx} 里声明，这里只给值。
%    \begin{macrocode}
\hit@presetup[lang=en]{ term = {
  chapter-name      = {Chapter\enskip,},
  appendixname      = {Appendix},
  contentsname      = {Contents},
  listfigurename    = {List of Figures},
  listtablename     = {List of Tables},
  figurename        = {Figure},
  tablename         = {Table},
  caption-algorithm-en = {Algorithm},
  equationname      = {Equation},
} }
%    \end{macrocode}
%
% 答辩决议表的表头。带 \verb|\\| 的几条是要在格子里换行的：评阅人那格三行，
% 「答辩委员会成员」与「委员」在纸质表上是竖排，逐字一行。
%    \begin{macrocode}
\hit@presetup[lang=en]{ term = {
  resolution-reviewer    = {Reviewers},
  resolution-name        = {Name},
  resolution-title       = {Title\\(Doctoral\\Supervisor?)},
  resolution-affiliation = {Affiliation},
  resolution-discipline  = {Discipline},
  resolution-role        = {Role},
  resolution-committee   = {\rotatebox{90}{Defense Committee}},
  resolution-chair       = {Chair},
  resolution-member      = {\rotatebox{90}{Member}},
  resolution-secretary   = {Secretary},
  resolution-text  = {Resolution of the defense committee (comments on the dissertation and whether the doctoral degree is recommended):},
} }
%    \end{macrocode}
%    \begin{macrocode}
\hit@presetup[lang=zh]{ term = {
  chapter-name      = {第,章},
  appendixname      = {附录},
  contentsname      = {目录},
  listfigurename    = {图清单},
  listtablename     = {表清单},
  figurename        = {图},
  tablename         = {表},
  caption-algorithm-zh = {算法},
  equationname      = {公式},
} }
%    \end{macrocode}
%
% 参考文献与索引各有两处要用：排出来的那条章标题跟文档语言走，英文目录里的条目
% 永远是英文。两半都写成成对的常量，\pkg{ctex} 的 \option{bibname} 与
% \option{indexname} 由它们接出来，不借 \pkg{ctex} 的键存中文那半。
%    \begin{macrocode}
\hit@presetup{ term = {
  bibliography-heading-zh = {参考文献},
  bibliography-heading-en = {References},
  index-heading-zh        = {索引},
  index-heading-en        = {Index},
} }
\hit@presetup[lang=zh]{ term = {
  bibname   = {<bibliography-heading-zh>},
  indexname = {<index-heading-zh>},
} }
\hit@presetup[lang=en]{ term = {
  bibname   = {<bibliography-heading-en>},
  indexname = {<index-heading-en>},
} }
\hit@presetup{ term = {
  list-of-figures-in-english-toc   = {List of figures},
  list-of-tables-in-english-toc    = {List of tables},
  list-of-equations-in-english-toc = {List of equations},
  listequationname  = {公式清单},
  abstract-heading-zh    = {摘要},
  abstract-in-english-toc-zh    = {Abstract (In Chinese)},
  abstract-heading-en    = {Abstract},
  abstract-in-english-toc-en    = {Abstract (In English)},
} }
\hit@presetup{ term = {
  abstract-keywords-zh = {关键词：},
  abstract-keywords-en = {Keywords:},
} }
\hit@presetup{ term = {
  abstract-keywords-separator-zh = {；},
  abstract-keywords-separator-en = {,\space},
} }
\ExplSyntaxOn

%    \end{macrocode}
%
% 双语题注里英文那条用的图表名。\cs{fnum@figure} 会在名字和编号之间放一个
% \cs{nobreakspace}，中文要这个空格、英文不要，所以默认值末尾带一个负细空格把它
% 抵掉。要换成别的写法直接覆盖这两个键。
%    \begin{macrocode}

\ExplSyntaxOff
\hit@presetup{ term = {
  caption-figure-prefix-en = {Fig.$\!$},
  caption-table-prefix-en  = {Table$\!$},
} }
\ExplSyntaxOn
\hit@preconfig{ \hit@setup@date }
%    \end{macrocode}
%
% 不写日期就取编译当天。这里存的是 \texttt{yyyy-mm-dd} 三段数字，中文封面还是英文
% 封面、排到年月还是年月日，由字段声明时的样式与精度定，见
% \file{src/engine/hit-keylist.dtx}。
%
% 外面套一层 \cs{use:x} 是要把今天固化成字面数字：留着 \cs{int\_use:N} 不展开，
% 解析那头拿到的是这几个记号本身，认不出是 ISO。
%
% 精度跟着值走（见 \file{src/engine/hit-date.dtx}），所以默认值只写到\emph{月}
% ——四处落点规范要的都是年月。用户写全 \texttt{yyyy-mm-dd} 就排到日。
% 今天取东八区，见 \file{src/engine/hit-engine.dtx}。
%    \begin{macrocode}
\use:x
  {
    \exp_not:N \hit@presetup
      {
        info =
          {
            date-zh = { \int_use:N \g__hit_engine_today_year_int -
              \int_use:N \g__hit_engine_today_month_int } ,
            date-en = { \int_use:N \g__hit_engine_today_year_int -
              \int_use:N \g__hit_engine_today_month_int } ,
          }
      }
  }
\ExplSyntaxOff
%    \end{macrocode}
%
% 封面上那行学生类别，学术与专业各有一条文本，各自是一个常量；元信息字段
% \option{degree-type} 只作覆盖用，不写就由查表按同名类选项挑一条。默认值不往
% 字段里写：模板给的默认值与用户自己写的值挤在同一个槽里，分不出这个字段是空着
% 还是被默认值填过，用户写了 \option{degree-type}\texttt{=professional} 会看不出变化。
%    \begin{macrocode}
\hit@presetup{ term = {
  degree-type-academic-zh     = {学术学位论文},
  degree-type-academic-en     = {Academic Degree},
  degree-type-professional-zh = {专业学位论文},
  degree-type-professional-en = {Professional Degree},
} }
\hit@presetup[degree-level=postdoc]{ term = {
  base-degree-level-zh           = {博士},
  base-degree-level-en           = {Doctor},
  titlepage-degree-attribute = {Doctoral},
} }
\ExplSyntaxOn

%    \end{macrocode}
%
% 添加博后选项
% \changes{v3.0.0}{2020/05/21}{添加博后选项}
% \changes{v3.1d}{2025/03/03}{统一将一堆 \cmd{\gdef} 放到了一起}
%    \begin{macrocode}

\ExplSyntaxOff
\hit@presetup[degree-level=doctor]{ term = {
  base-degree-level-zh           = {博士},
  base-degree-level-en           = {Doctor},
  titlepage-degree-attribute = {Doctoral},
} }
\hit@presetup[degree-level=master]{ term = {
  base-degree-level-zh           = {硕士},
  base-degree-level-en           = {Master},
  titlepage-degree-attribute = {Master's},
} }
\ExplSyntaxOn
\ExplSyntaxOff
\ExplSyntaxOn
%    \end{macrocode}
%
% 内封上这两处按校区与学位挑的是印在纸上的字，不是版面结构，所以归词条表而不
% 写进排版代码；\cs{hitsetup} 里写显式值就能覆盖。
%
% 深圳硕士那一档的英文句子里不带学位属性词（写死的 \texttt{Master}），
% 与本部那一档不是同一句，所以单列一条，不走 \option{titlepage-degree-attribute}。
%
% 两条名字都不带语言后缀。内封的中文半页与英文半页各排各的，
% \option{titlepage-degree-line} 只有中文、\option{titlepage-degree-sentence}
% 只有英文，两者不是同一句话的两个语种，凑不成一对。带后缀就意味着分语言、
% 必须有兄弟，见 \file{scripts/check-unused-const.py}。
%    \begin{macrocode}
\ExplSyntaxOff
\hit@presetup{ term = {
  titlepage-degree-sentence = {Dissertation for the <titlepage-degree-attribute> Degree in <discipline-en>},
  titlepage-degree-line     = {<base-degree-of-discipline-zh>},
} }
\hit@presetup[campus=harbin,degree-level=master ; campus=harbin,degree-level=doctor]{ term = {
  titlepage-degree-sentence = {Dissertation for the <titlepage-degree-attribute> Degree},
} }
\hit@presetup[campus=shenzhen,degree-level=master]{ term = {
  titlepage-degree-sentence = {Dissertation for the Master Degree},
} }
\hit@presetup[campus=harbin,degree-level=master ; campus=harbin,degree-level=doctor ; campus=shenzhen,degree-level=master]{ term = {
  titlepage-degree-line = {<base-degree-level-zh>},
} }
\hit@presetup{ term = {
  base-degree-of-discipline-zh               = {<discipline-zh><base-degree-level-zh>},
  base-degree-of-discipline-en               = {<base-degree-level-en> of <discipline-en>},
  titlepage-author-zh                        = {<base-degree-level-zh>研究生},
  cover-classified-index-postdoc      = {分类号},
  base-secrecy-postdoc          = {密级},
  cover-udc-postdoc                      = {\hit@spread@by{0.4\ccwd}{UDC}},
  base-stage-proposal                                = {开题},
  base-stage-interim                                = {中期},
  base-stage-document-type                          = {报告},
  cover-number-postdoc                = {编号},
  base-report-postdoc                            = {博士后研究工作报告},
  cover-work-finish-date-postdoc      = {工作完成日期},
  cover-report-submit-date-postdoc    = {报告提交日期},
  titlepage-author-postdoc-zh                = {\hit@spread@to{7em}{博士后姓名}},
  titlepage-co-supervisor-postdoc-zh         = {\hit@spread@to{7em}{合作导师}},
  titlepage-station-postdoc               = {\hit@spread@to{7em}{流动站名称}},
  titlepage-major-postdoc-zh             = {\hit@spread@to{7em}{专业名称}},
  titlepage-research-start-date-postdoc   = {研究工作起始时间},
  titlepage-research-end-date-postdoc     = {研究工作期满时间},
  titlepage-separator-postdoc                            = {：},
  base-document-type-level-bachelor                     = {本科},
  base-document-type-bachelor                    = {毕业论文（设计）},
  titlepage-affiliation-bachelor-zh          = {学院 },
  titlepage-student-id-bachelor           = {学号},
  titlepage-major-bachelor-zh                = {专业},
  titlepage-supervisor-bachelor-zh           = {指导教师},
  cover-title-label-bachelor               = {题目},
} }
\hit@presetup[degree-level=bachelor]{ term = {
  base-degree-level-zh                     = {学士 },
  titlepage-author-zh               = {本科生},
} }
\ExplSyntaxOn
%    \end{macrocode}
%
% 更改博后封皮中字间距
% \changes{v3.0.12}{2020/11/02}{更改博后封皮中字间距}
%
% 添加博后选项
% \changes{v3.0.0}{2020/05/21}{添加博后选项}
% \changes{v3.1d}{2025/03/03}{删了一些没用的定义，更改哈尔滨本科的论文名。\issue[172]}
%
% 添加本科开题封皮中title定义
% \changes{v3.0.0}{2020/05/21}{添加本科开题封皮中title定义}
% \changes{v3.1d}{2025/03/03}{哈尔滨本科由院系改成了学院。\issue[172]}
%
% 本科没有“学位简称加士”这套说法，几个词都要单独给。三个校区已经统一，
% 毕业论文与院系的叫法都只剩一份。\texttt{bachelor-affiliation-field-label-weihai}
% 一起删了，排版代码从来没用过它。
%
% \changes{v3.1b}{2022/06/06}{修改深圳校区本科生第二封面中“学生”为“姓名”}
%    \begin{macrocode}
\ExplSyntaxOff
\hit@presetup{ term = {
  titlepage-student-bachelor = {学生},
  base-document-type                = {学位论文},
  titlepage-date-bachelor-zh    = {\hit@spread@to{4em}{日期}},
} }
\hit@presetup{ term = {
  cover-teacher-comment-bachelor = {指导教师评语：},
} }
\hit@presetup{ term = {
  cover-teacher-signature-bachelor = {指导教师签字：},
} }
\hit@presetup{ term = {
  cover-check-date-bachelor = {检查日期：},
} }
\ExplSyntaxOn
\ExplSyntaxOff
\hit@presetup{ term = {
  cover-title-prefix                    = {\hit@spread@to{3em}{题目}},
  titlepage-affiliation       = {院\hspace{3\ccwd}（系）},
  titlepage-major             = {学\hspace{4\ccwd}科},
  titlepage-supervisor        = {\hit@spread@to{6em}{导师}},
  titlepage-student           = {\hit@spread@to{6em}{研究生}},
  titlepage-student-id        = {\hit@spread@to{6em}{学号}},
  titlepage-date              = {\hit@if@option@TF{proposal}{<base-stage-proposal>}%
{\hit@if@option@T{interim}{<base-stage-interim>}}<base-stage-document-type>日期},
  cover-enroll-date       = {\hit@spread@to{6em}{入学时间}},
  cover-title-label            = {\hit@spread@to{6em}{论文题目}},
  titlepage-affiliation-major = {\hit@spread@to{6em}{学科专业}},
  base-university-zh                                 = {哈尔滨工业大学},
  titlepage-institution-zh                     = {授予学位单位},
  titlepage-date-zh                       = {答辩日期},
  titlepage-affiliation-zh                = {所在单位},
} }
\ExplSyntaxOn
%    \end{macrocode}
%
% \cs{hit@stage@proposal} 这条留在值里判断，不能提到覆盖表。\option{stage} 是报告
% 的选项，\file{hit-thesis.cls} 里没有 \cs{g\_\_hit\_interim\_bool}，提到外面
% 加载期就会报 \texttt{Missing number}。留在值里则只有真的排到这一项时才求值，
% 而 book 从不用它。
%
% 添加哈尔滨校区选项
% \changes{v3.0.0}{2020/05/21}{添加哈尔滨校区选项}
% \changes{v3.1d}{2025/03/03}{哈尔滨本科多了一个“学校”。\issue[172]}
%
%    \begin{macrocode}
\ExplSyntaxOff
\hit@presetup[degree-level=bachelor]{ term = {
  titlepage-institution-zh = {学校},
} }
\ExplSyntaxOn
\ExplSyntaxOff
\hit@presetup{ term = {
  titlepage-major-zh            = {学科},
  titlepage-degree-applied-zh               = {申请学位},
  titlepage-supervisor-zh           = {导师},
  titlepage-associate-supervisor-zh = {副导师},
  titlepage-co-supervisor-zh        = {合作导师},
  base-title-separator-zh                  = {：},
  titlepage-author-en               = {Candidate},
  titlepage-supervisor-en           = {Supervisor},
  titlepage-associate-supervisor-en = {Associate Supervisor},
  titlepage-co-supervisor-en        = {Co-Supervisor off Campus},
  titlepage-degree-applied-en               = {Academic Degree Applied for},
  titlepage-major-en            = {Specialty},
  titlepage-affiliation-en          = {Affiliation},
} }
\ExplSyntaxOn
%    \end{macrocode}
%
% \changes{v3.1b}{2022/06/06}{本部研究生第二封面按照书写指南进行修改}
% 本部研究生第二封面按照书写指南进行修改。学科”改为“学科或类别”。
% 虽然只是深圳的要求，但是懒得写判断了，我估计用到的人不多，还有我怕明年深圳再改一波又白写了。
% \changes{v3.1c}{2023/04/16}{根据 \issue[97] 修改合作导师的名称}
%
% \changes{v3.1c}{2023/04/16}{根据 \issue[97] 修改合作导师的英文}
%
%    \begin{macrocode}
\ExplSyntaxOff
\hit@presetup[campus=harbin,degree-level=master]{ term = {
  titlepage-major-zh = {学科或类别},
} }
\ExplSyntaxOn

%    \end{macrocode}
%
% \changes{v3.1b}{2022/06/06}{修改深圳硕士英文封面的 Defence 为 Defense.\issue[97]}
% 深圳硕士英文封面是 Defense，其他为 Defence.
%    \begin{macrocode}

\ExplSyntaxOff
\hit@presetup{ term = {
  titlepage-date-en                           = {Date of Defence},
  titlepage-institution-en                         = {Degree-Conferring-Institution},
  base-university-en                                     = {Harbin Institute of Technology},
  base-title-separator-en                            = {:},
  titlepage-classified-index-national-zh      = {国内图书分类号},
  titlepage-classified-index-international-zh = {国际图书分类号},
  titlepage-classified-index-national-en      = {Classified Index},
  titlepage-classified-index-international-en = {U.D.C},
  table-of-contents-page-en                           = {Page},
  table-of-contents-table-en                          = {Table},
  table-of-contents-figure-en                         = {Figure},
  table-of-contents-undefined-en                      = {(Error! English heading not defined.)},
  base-secrecy                   = {密级},
  titlepage-school-code                      = {学校代码},
  base-school-code                                     = {10213},
} }
\ExplSyntaxOn
%    \end{macrocode}
%
% 终稿独有那几格的回落值，同样照范例的占位形。学科给两个方框，
% 「学位」那一行经 \option{base-degree-of-discipline} 拼出来就是范例内封上的
% \texttt{□□博士}；英文学科四个方框，拼出 \texttt{Doctor~of~□□□□}。
% 密级范例印「公开」，两个分类号印四个叉——叉是宋体形，与范例同（\texttt{×}
% 平时归半角标点类取西文字形，这里逐字用 \cs{hit@engine@cjk@char} 点名按汉字
% 排）。范例的占位形只到这里——副导师、
% 联合导师范例上是空行，照旧空着。
%    \begin{macrocode}
\ExplSyntaxOff
\hit@presetup{ info = {
  title-en = {RESEARCH ON KEY TECHNOLOGIES OF PARTIAL POROUS EXTERNALLY PRESSURIZED GAS BEARING},
  author-en     = {□□□},
  supervisor-en = {Prof.□□□},
  affiliation-en = {School of Mechatronics Engineering},
  major-en      = {Mechanical Manufacturing and Automation},
  discipline-zh = {□□},
  discipline-en = {□□□□},
  secret-level  = {公开},
  classified-index-national      = {\hit@engine@cjk@char{×}\hit@engine@cjk@char{×}\hit@engine@cjk@char{×}\hit@engine@cjk@char{×}},
  classified-index-international = {\hit@engine@cjk@char{×}\hit@engine@cjk@char{×}\hit@engine@cjk@char{×}\hit@engine@cjk@char{×}},
} }
\ExplSyntaxOn

%    \end{macrocode}
%
%    \begin{macrocode}

\ExplSyntaxOff
\hit@presetup[campus=shenzhen,degree-level=master]{ term = {
  titlepage-date-en = {Date of Defense},
} }
\hit@presetup{ term = {
  base-heading-spread             = {1em},
  conclusion-heading-zh      = {结论},
  conclusion-heading-en      = {Conclusions},
  acknowledgements-heading-zh = {致谢},
  acknowledgements-heading-en = {Acknowledgements},
  resume-heading-zh          = {个人简历},
} }
\ExplSyntaxOn

%    \end{macrocode}
%
% 博后这三个标题与别的学位不同：正文标题不加字距，简历要写明是博士后的。
%    \begin{macrocode}

\ExplSyntaxOff
\hit@presetup[degree-level=postdoc]{ term = {
  base-heading-spread    = {},
  resume-heading-zh = {博士后个人简历},
} }
\ExplSyntaxOn
\ExplSyntaxOff
\hit@presetup{ term = {
  resume-corresponding-address-zh = {永久通讯地址},
  resume-corresponding-address-en = {Corresponding Address},
  resume-heading-en                = {Resume},
  declarations-heading-zh          = {哈尔滨工业大学学位论文原创性声明和使用权限},
  declarations-heading-en          = {Statement of copyright and Letter of authorization},
  resolution-heading-zh            = {学位论文评阅人、答辩委员会名单及答辩决议},
  resolution-heading-en            = {List of Dissertation Reviewers and Defense Committee and Defense Resolution},
  resolution-reviewer        = {评阅人\\（根据实际\\人数填写）},
  resolution-name            = {姓名},
  resolution-title           = {职称（是否博导）},
  resolution-affiliation     = {工作单位},
  resolution-discipline      = {所在学科},
  resolution-role            = {职务},
  resolution-committee       = {答\\辩\\委\\员\\会\\成\\员},
  resolution-chair           = {主席},
  resolution-member          = {委\\员},
  resolution-secretary       = {秘书},
  resolution-text      = {答辩委员会决议（对论文的评语及是否建议授予博士学位等）：},
} }
\ExplSyntaxOn
%    \end{macrocode}
%
% 添加博后封皮title的选项
% \changes{v3.0.0}{2020/05/21}{添加博后封皮title的选项}
%    \begin{macrocode}


%    \end{macrocode}
%
%    \begin{macrocode}
\ExplSyntaxOff
\hit@presetup{ term = {
  declarations-author-signature = {作者签名：},
} }
\hit@presetup{ term = {
  declarations-supervisor-signature = {导师签名：},
} }
\hit@presetup{ term = {
  declarations-date = {日期：},
} }
\hit@presetup{ term = {
  list-of-symbols-heading-zh = {物理量名称及符号表},
  list-of-abbreviations-heading-zh = {缩略语表},
  list-of-symbols-and-abbreviations-heading-zh = {符号及缩略语},
} }
\hit@presetup{ term = {
  list-of-symbols-heading-en = {List of physical quantities and symbols},
  list-of-abbreviations-heading-en = {List of abbreviations},
  list-of-symbols-and-abbreviations-heading-en = {List of symbols and abbreviations},
} }
\hit@presetup{ term = {
  mainmatter-eqdenote-label-zh = {式中},
} }
\hit@presetup{ term = {
  mainmatter-eqdenote-label-en = {where},
} }
\hit@presetup{ term = {
  mainmatter-eqdenote-dash = {——},
} }
\hit@presetup{ term = {
  declarations-authorization = {学位论文使用权限},
} }
%    \end{macrocode}
%
% \changes{v3.0.22}{2022/05/16}{深圳校区中期模板封面范例中没有“哈工大（深圳）制”字样。\issue[126]}
% 深圳校区的封面底部没有「哈工大（深圳）制」这一行。另外三个校区学位档次的底部标记
% 都在下面，唯独深圳这个不排，2022 年按 \issue[126] 去掉的，别再补回来。
%    \begin{macrocode}
\hit@presetup{ term = {
  cover-school-bottom-mark-harbin-master = { 
   \begin{center} 
     \songti\sanhao\textbf{研究生院制} 
   \end{center} 
   \begin{center} 
     \songti\sanhao\textbf{二〇一四年九月} 
   \end{center} 
},
  cover-school-bottom-mark-harbin-doctor = { 
   \begin{center} 
     \songti\sanhao\textbf{研究生院制} 
   \end{center} 
   \begin{center} 
     \songti\sanhao\textbf{二〇一四年九月} 
   \end{center} 
},
} }
\ExplSyntaxOn

%    \end{macrocode}
%
% \changes{v3.1c}{2023/04/16}{根据 \issue[151] 将本部本科开题下方的标识更换为 \texttt{.eps} 文件}
%    \begin{macrocode}

\ExplSyntaxOff
\hit@presetup{ term = {
  cover-school-bottom-mark-bachelor = { 
   \begin{center} 
     \hit@cover@bottommark 
   \end{center} 
},
} }
\hit@presetup{ term = {
  cover-note-interim-shenzhen-doctor = {
\thispagestyle{empty}
{\begin{center}
\songti\sanhao\textbf{\hit@spread@to{6em}{说明}}
\end{center}
}
\vspace{1cm}\songti\hit@fontsize@x{sihao}{2}
\begin{enumerate}[itemsep=-4bp]
\item 博士研究生学位论文中期报告一般在入学后的第五学期末完成。
\item 此报告由学生本人填写，请导师签字后提交一份给检查小组。答辩结束后由检查组长签署意见。
\item 报告经中期报告检查组长签字同意后，由学科部保存，以备论文答辩时参考。研究生院学生培养处将对研究生的学位论文中期报告进行抽查。
\item 报告统一用A4纸打印。
\end{enumerate}
},
} }
\ExplSyntaxOn

%    \end{macrocode}
%
% 添加深圳中期报告内容
% \changes{v3.0.0}{2020/05/21}{添加深圳中期报告内容}
%    \begin{macrocode}

\ExplSyntaxOff
\hit@presetup{ term = {
  declarations-authorization-text = {%
学位论文是研究生在哈尔滨工业大学攻读学位期间完成的成果，知识产权归属哈尔滨工业大学。学位论文的使用权限如下： 
\par
（1）学校可以采用影印、缩印或其他复制手段保存研究生上交的学位论文，并向国家图书馆报送学位论文；（2）学校可以将学位论文部分或全部内容编入有关数据库进行检索和提供相应阅览服务；（3）研究生毕业后发表与此学位论文研究成果相关的学术论文和其他成果时，应征得导师同意，且第一署名单位为哈尔滨工业大学。 
\par
保密论文在保密期内遵守有关保密规定，解密后适用于此使用权限规定。 
\par
本人知悉学位论文的使用权限，并将遵守有关规定。},
} }
\ExplSyntaxOn

%    \end{macrocode}
%
% 添加威海校区原创性声明内容
% \changes{v3.0.16}{2021/11/29}{添加威海校区原创性声明内容}
% \changes{v3.1d}{2025/03/03}{哈尔滨本科的原创性声明的标题也改了。\issue[172]}
%    \begin{macrocode}

%    \end{macrocode}
%
% \changes{v3.1b}{2022/06/06}{更改深圳校区本科生原创性声明标题}
% 三个校区已经统一，标题只剩这一个，分成两行排。
%    \begin{macrocode}

\ExplSyntaxOff
\hit@presetup{ term = {
  declarations-heading-line-one-bachelor = {哈尔滨工业大学本科毕业论文（设计）},
  declarations-heading-line-two-bachelor = {原创性声明和使用权限},
  declarations-heading-bachelor-zh = {<declarations-heading-line-one-bachelor><declarations-heading-line-two-bachelor> },
  declarations-heading-short-bachelor = {原创性声明},
  declarations-statement-bachelor = {本科毕业论文（设计）原创性声明},
  declarations-authorization-bachelor = {本科毕业论文（设计）使用权限},
} }
\ExplSyntaxOn
%    \end{macrocode}
%    \begin{macro}{\hit@term@declarations@statement@text@bachelor}
% \changes{v3.1d}{2025/03/03}{修改哈尔滨本科原创性声明的内容。\issue[172]}
%
% 三个校区已经统一，正文只剩这一块。末尾那句 AI 工具是新规范加的，
% 与前文同一段，不另起一段。
%    \begin{macrocode}
\ExplSyntaxOff
\hit@presetup{ term = {
  declarations-statement-text-bachelor = {%
    本人郑重声明：此处所提交的本科毕业论文（设计）《<title-flat-zh>》，是本人在导师指导下，在哈尔滨工业大学攻读学士学位期间独立进行研究工作所取得的成果，且毕业论文（设计）中除已标注引用文献的部分外不包含他人完成或已发表的研究成果。对本毕业论文（设计）的研究工作做出重要贡献的个人和集体，均已在文中以明确方式注明。本毕业论文（设计）对使用AI工具的情况进行了明确标注。
  },
} }
\ExplSyntaxOn

%    \end{macrocode}
%    \end{macro}
%    \begin{macro}{\hit@term@declarations@authorization@text@bachelor}
% \changes{v3.1d}{2025/03/03}{修改哈尔滨本科知识产权的内容。\issue[172]}
%    \begin{macrocode}

\ExplSyntaxOff
\hit@presetup{ term = {
  declarations-authorization-text-bachelor = {%
  本科毕业论文（设计）是本科生在哈尔滨工业大学攻读学士学位期间完成的成果，知识产权归属哈尔滨工业大学。本科毕业论文（设计）的使用权限如下： 
\par
  （1）学校可以采用影印、缩印或其他复制手段保存本科生上交的毕业论文（设计），并向有关部门报送本科毕业论文（设计）；（2）根据需要，学校可以将本科毕业论文（设计）部分或全部内容编入有关数据库进行检索和提供相应阅览服务；（3）本科生毕业后发表与此毕业论文（设计）研究成果相关的学术论文和其他成果时，应征得导师同意，且第一署名单位为哈尔滨工业大学。 
\par
  保密论文在保密期内遵守有关保密规定，解密后适用于此使用权限规定。 
\par
  本人知悉本科毕业论文（设计）的使用权限，并将遵守有关规定。 
},
} }
\ExplSyntaxOn

%    \end{macrocode}
%    \end{macro}
%    \begin{macro}{\hit@term@declarations@statement}
%    \begin{macrocode}

\ExplSyntaxOff
\hit@presetup{ term = {
  declarations-statement = {学位论文原创性声明},
} }
\ExplSyntaxOn

%    \end{macrocode}
%    \end{macro}
%    \begin{macro}{\hit@term@declarations@statement@text}
%    \begin{macrocode}

\ExplSyntaxOff
\hit@presetup{ term = {
  declarations-statement-text = {%
本人郑重声明：此处所提交的学位论文《<title-flat-zh>》，是本人在导师指导下，在哈尔滨工业大学攻读学位期间独立进行研究工作所取得的成果，且学位论文中除已标注引用文献的部分外不包含他人完成或已发表的研究成果。对本学位论文的研究工作做出重要贡献的个人和集体，均已在文中以明确方式注明。},
} }
\ExplSyntaxOn

%    \end{macrocode}
%    \end{macro}
%
% \changes{v3.1b}{2022/06/06}{修改深圳校区硕士创新性成果的标题。\issue[97]}
% \changes{v3.1d}{2025/03/03}{修改哈尔滨校区本科创新性成果的标题。\issue[220]}
% 成果附录的标题。
% 理工与人文社科两份研究生指南（2.17／（十七））、本科指南、博士范例、本科
% 范例一个叫法：“攻读□士学位期间取得创新性成果”；“发表的论文及其他成果”是
% 老规范的字样，不用。博后和英文版另有叫法，按覆盖顺序排在后面。
%    \begin{macrocode}

\ExplSyntaxOff
\hit@presetup{ term = {
  publications-heading-zh = {攻读<base-degree-level-zh>学位期间取得创新性成果},
} }
\hit@presetup[degree-level=postdoc]{ term = {
  publications-heading-zh = {博士后期间的研究成果},
} }
\hit@presetup[lang=en]{ term = {
  publications-heading-zh = {Author's Publications},
} }
\hit@presetup{ term = {
  publications-doctoral-heading-zh = {博士生期间的研究成果},
  publications-doctoral-heading-en = {Publications During Doctoral Study},
} }
\ExplSyntaxOn

%    \end{macrocode}
%
% 博后和英文版附录title
% \changes{v3.0.0}{2020/05/21}{博后和英文版附录title}
%    \begin{macrocode}

\ExplSyntaxOff
\hit@presetup{ term = {
  declarations-date-blank = {\hspace{2.5em}年\hspace{1.5em}月\hspace{1.5em}日},
} }
\ExplSyntaxOn

%    \end{macrocode}
%
% \changes{v3.1b}{2022/06/06}{修正英文目录中的硕士期间发表文章。\issue[97]}
%    \begin{macrocode}

\ExplSyntaxOff
\hit@presetup{ term = {
  publications-heading-en = {Innovative achievements for Ph.D},
} }
\ExplSyntaxOn
\ExplSyntaxOff
\hit@presetup[degree-level=master]{ term = {
  publications-heading-en = {Innovative achievements for Master},
} }
\ExplSyntaxOn
\ExplSyntaxOff
\hit@presetup{ term = {
  base-hi               = {嗨！thesis},
  base-brace-left-zh    = {（},
  base-brace-right-zh   = {）},
  base-brace-left-en    = {(},
  base-brace-right-en   = {)},
} }
\ExplSyntaxOn
%    \end{macrocode}
%
%    \begin{macrocode}
\hit@presetup@scope{stage=!final}
%    \end{macrocode}
%
% \emph{深圳本科开题与中期首页的字。} 取自 iota-hit 的
% \texttt{src/config/terms.typ}，那边是逐字量 2026 年 9 月版两份原件
% （\file{深圳本科毕业论文（设计）开题报告.doc} 与同版中期）得来的。本仓库没有
% 原件，照抄不复核。
%
% \emph{校名印的是第三个字面。} 首页与页眉都印「哈尔滨工业大学深圳校区」，
% 而研究生的页眉印「哈尔滨工业大学（深圳）」，两回事。
%
% \emph{文种不带「本科」两个字。} 本部那两份印「本科毕业论文（设计）开题报告」，
% 深圳这两份印「毕业论文（设计）开题报告」。2026 年 9 月版不再按论文还是设计
% 换字，改由首页第一行那对勾选框表示。
%
% \emph{六格的标签自带全角冒号}，英文版是「中文／ English：」的双语形式。
%    \begin{macrocode}
\ExplSyntaxOff
\hit@presetup[campus=shenzhen,degree-level=bachelor]{ term = {
  cover-university-shenzhen-bachelor-zh      = {哈尔滨工业大学深圳校区},
  cover-university-shenzhen-bachelor-en      = {Harbin Institute of Technology, Shenzhen},
  cover-form-line-shenzhen-bachelor-zh       = {毕业论文（设计）},
  cover-form-line-shenzhen-bachelor-en       = {Undergraduate Thesis (Design)},
  cover-title-label-shenzhen-bachelor-zh     = {题目：},
  cover-title-label-shenzhen-bachelor-en     = {题目/ Title:},
  cover-author-label-shenzhen-bachelor-zh    = {姓名：},
  cover-author-label-shenzhen-bachelor-en    = {姓名/ Name：},
  cover-student-id-label-shenzhen-bachelor-zh = {学号：},
  cover-student-id-label-shenzhen-bachelor-en = {学号/ Student~ID：},
  cover-affiliation-label-shenzhen-bachelor-zh = {学院：},
  cover-affiliation-label-shenzhen-bachelor-en = {学院/ School：},
  cover-speciality-label-shenzhen-bachelor-zh = {专业：},
  cover-speciality-label-shenzhen-bachelor-en = {专业/ Major：},
  cover-supervisor-label-shenzhen-bachelor-zh = {指导教师：},
  cover-supervisor-label-shenzhen-bachelor-en = {指导教师/ Supervisor：},
  cover-date-label-shenzhen-bachelor-zh      = {日期：},
  cover-date-label-shenzhen-bachelor-en      = {日期/ Date：},
} }
\hit@presetup[campus=shenzhen,degree-level=bachelor,stage=proposal]{ term = {
  cover-stage-line-shenzhen-bachelor-zh = {开题报告},
  cover-stage-line-shenzhen-bachelor-en = {Proposal~Report},
} }
\hit@presetup[campus=shenzhen,degree-level=bachelor,stage=interim]{ term = {
  cover-stage-line-shenzhen-bachelor-zh = {中期报告},
  cover-stage-line-shenzhen-bachelor-en = {Interim~Report},
} }
\ExplSyntaxOn
%    \end{macrocode}
%
%    \begin{macrocode}
\ExplSyntaxOff
%    \end{macrocode}
%
% 报告的页眉整句也是常量，取值按变体给，与学位论文那边同一种做法，用户改得了。
%
% 深圳硕士的前置页面与正文页眉不一样，所以有两个键。其余变体只用
% \option{report-header}。
%    \begin{macrocode}
\hit@presetup[campus=weihai,degree-level=master]{ term = {
  header-report = {<base-university-zh>（<base-campus>）<base-degree-level-zh>学位论文开题报告},
} }
\hit@presetup[campus=shenzhen,degree-level=master,stage=proposal]{ term = {
  header-report = {<base-university-zh>（<base-campus>）<base-degree-level-zh>学位论文开题报告},
} }
\hit@presetup[campus=shenzhen,degree-level=master,stage=interim]{ term = {
  header-report = {<base-university-zh>（<base-campus>）<base-stage-interim>报告},
} }
\hit@presetup[campus=shenzhen,degree-level=master]{ term = {
  header-report-front = {<base-university-zh>（<base-campus>）<base-degree-level-zh>学位论文开题报告},
} }
%% 深圳本科那两份的页眉与首页文种那一行同字：校名加表单名。校名是第三个字面
%% （「哈尔滨工业大学深圳校区」），与研究生页眉的「哈尔滨工业大学（深圳）」不是
%% 一回事，所以从首页那三个词条拼，不从 base 拼。
\hit@presetup[campus=shenzhen,degree-level=bachelor,stage=proposal]{ term = {
  header-report = {<cover-university-shenzhen-bachelor-zh><cover-form-line-shenzhen-bachelor-zh><cover-stage-line-shenzhen-bachelor-zh>},
} }
\hit@presetup[campus=shenzhen,degree-level=bachelor,stage=interim]{ term = {
  header-report = {<cover-university-shenzhen-bachelor-zh><cover-form-line-shenzhen-bachelor-zh><cover-stage-line-shenzhen-bachelor-zh>},
} }
%% 英文版那两份的页眉只有表单名，没有校名——逐字量 iota-hit 对原件的复刻件
%% （英文版开题、中期的第二页起）：印的就是「Undergraduate Thesis (Design)
%% Proposal Report」，居中，小五，一个校名都没有。研究生那两份英文版才是
%% 「表单名 ＋ 校名」，那一档这里还没有原件。
\hit@presetup[campus=shenzhen,degree-level=bachelor,stage=proposal,lang=en]{ term = {
  header-report = {<cover-form-line-shenzhen-bachelor-en>~<cover-stage-line-shenzhen-bachelor-en>},
} }
\hit@presetup[campus=shenzhen,degree-level=bachelor,stage=interim,lang=en]{ term = {
  header-report = {<cover-form-line-shenzhen-bachelor-en>~<cover-stage-line-shenzhen-bachelor-en>},
} }
%% 英文版那两份的参考文献标题印「References」，只给中文那一条英文报告会报
%% 「词条没取值」。
\hit@presetup{ term = {
  bibliography-heading-zh = {参考文献},
  bibliography-heading-en = {References},
} }
%% 报告这一档三个名字也按语种分开，英文版印 Contents／Figure／Table（iota-hit
%% 对原件的复刻件里题注是「Figure 1-1」「Table 1-1」）。
\hit@presetup[lang=zh]{ term = {
  contentsname = {\hit@spread@by{1em}{目录}},
  figurename   = {图},
  tablename    = {表},
} }
\hit@presetup[lang=en]{ term = {
  contentsname = {Contents},
  figurename   = {Figure},
  tablename    = {Table},
} }
\ExplSyntaxOn
\hit@preconfig{ \hit@setup@date }
%    \end{macrocode}
%
% \changes{v3.1d}{2025/03/03}{根据 \issue[220] 去掉本部答辩日期中的“日”}
%
% 不写日期就取编译当天（东八区），写成 ISO 交给格式化那一层；排到哪一级由
% 值自己定。哈尔滨与威海到年月；深圳本科的开题中期封面要年月日
% （\issue{220}），默认值单独给全。
%    \begin{macrocode}
\use:x
  {
    \exp_not:N \hit@presetup
      { info = { date-zh = { \int_use:N \g__hit_engine_today_year_int -
              \int_use:N \g__hit_engine_today_month_int } } }
  }
\use:x
  {
    \exp_not:N \hit@presetup [ campus = shenzhen , degree-level = bachelor ]
      { info = { date-zh = { \int_use:N \g__hit_engine_today_year_int -
              \int_use:N \g__hit_engine_today_month_int -
              \int_use:N \g__hit_engine_today_day_int } } }
  }
%    \end{macrocode}
%
% \changes{v3.1d}{2025/03/03}{删了一些没用的定义，更改哈尔滨本科的论文名。\issue[172]}
% 咱也不知道谁想的，其他校区的本科毕设叫“毕业设计（论文）”，哈尔滨的叫“毕业论文（设计）”
%
% \changes{v3.1d}{2025/03/03}{之前一口气删多了，恢复一下}
%
% \changes{v3.1c}{2023/04/16}{修改深圳硕士开题封面“学科”}
%
% 上面那张表只放默认值，下面按条件覆盖。后设的盖先设的，所以条件从宽到窄排。
%    \begin{macrocode}
\ExplSyntaxOff
\hit@presetup[degree-level=doctor]{ term = {
  base-degree-level-zh                = {博士},
  base-degree-of-discipline-zh = {<discipline-zh><base-degree-level-zh>},
  titlepage-author-zh          = {<base-degree-level-zh>研究生},
} }
\hit@presetup[degree-level=master]{ term = {
  base-degree-level-zh                = {硕士},
  base-degree-of-discipline-zh = {<discipline-zh><base-degree-level-zh>},
  titlepage-author-zh          = {<base-degree-level-zh>研究生},
} }
\hit@presetup[degree-level=bachelor]{ term = {
  base-degree-level-zh = {学士},
} }
\hit@presetup{ term = {
  caption-figure-prefix-en                              = {Fig.$\!$},
  caption-table-prefix-en                               = {Table$\!$},
  base-stage-proposal                                 = {开题},
  base-stage-interim                                 = {中期},
  base-stage-document-type                           = {报告},
  base-document-type-level-bachelor                      = {本科},
  base-document-type-bachelor                     = {毕业设计（论文）},
  cover-class-bachelor                 = {班号},
  cover-major-bachelor                 = {\hit@spread@to{4em}{专业}},
  cover-student-bachelor               = {\hit@spread@to{4em}{学生}},
  cover-student-id-bachelor            = {\hit@spread@to{4em}{学号}},
  cover-supervisor-bachelor            = {指导教师},
  cover-report-date-bachelor                  = {\hit@spread@to{4em}{日期}},
  cover-teacher-comment-bachelor       = {指导教师评语：},
  cover-teacher-signature-bachelor     = {指导教师签字：},
  cover-check-date-bachelor            = {检查日期：},
  cover-title-prefix                        = {\hit@spread@to{3em}{题目}},
  cover-affiliation           = {\hit@spread@to{6\ccwd}{院（系}）},
  cover-major                 = {\hit@spread@to{6em}{学科}},
  cover-supervisor            = {\hit@spread@to{6em}{导师}},
  cover-student               = {\hit@spread@to{6em}{研究生}},
  cover-student-id            = {\hit@spread@to{6em}{学号}},
  cover-report-date                  = {\hit@term@base@stage@proposal\hit@term@base@stage@document@type 日期},
  cover-enroll-date           = {\hit@spread@to{6em}{入学时间}},
  cover-title-label                = {\hit@spread@to{6em}{论文题目}},
  cover-affiliation-major     = {\hit@spread@to{6em}{学科专业}},
  base-university-zh                                     = {哈尔滨工业大学},
  base-document-type                              = {学位论文},
  cover-check-report-title                         = {\hit@spread@to{13.375em}{博士学位论文中期检查报告}},
  titlepage-institution-zh                         = {授予学位单位},
  titlepage-degree-applied-zh                         = {申请学位},
  titlepage-associate-supervisor-zh           = {副导师},
  titlepage-co-supervisor-zh                  = {联合导师},
  base-title-separator-zh                            = {：},
  titlepage-classified-index-national-zh      = {国内图书分类号},
  titlepage-classified-index-international-zh = {国际图书分类号},
  base-secrecy                   = {密级},
  titlepage-school-code                      = {学校代码},
  base-school-code                                     = {10213},
} }
%    \end{macrocode}
%
%    \begin{macrocode}
\hit@presetup[campus=harbin|weihai]{ term = {
  base-document-type-bachelor = {毕业论文（设计）},
} }
\hit@presetup[campus=harbin,degree-level=bachelor]{ term = {
  titlepage-institution-zh = {学校},
} }
\hit@presetup[stage=interim]{ term = {
  cover-report-date = {\hit@term@base@stage@interim\hit@term@base@stage@document@type 日期},
} }
\ExplSyntaxOn
\ExplSyntaxOff
\hit@presetup[campus=shenzhen,degree-level=master,stage=proposal]{ term = {
  cover-major = {\hit@spread@to{6\ccwd}{学科\,/\,专业}},
} }
\hit@presetup[campus=harbin,degree-level=master]{ term = {
  cover-affiliation = {\hit@spread@to{6\ccwd}{学院（部}）},
  % 这一处是压不是拉，九个字进八个字宽。\hit@spread@to 底层是 minipage 会折行，
  % 压不了，所以留 \makebox[s]。别处的分散对齐一律走 \hit@spread@to。
  cover-major       = {\makebox[8\ccwd][s]{学科/专业学位类别}},
} }
\hit@presetup[campus=harbin,degree-level=master,stage=proposal]{ term = {
  base-document-type = {学位},
} }
\ExplSyntaxOn
%    \end{macrocode}
%
% 报告封面上这四处按变体分岔的是字距字号，是排版取值不是结构，写进排版代码用户
% 盖不掉，所以做成词条。\option{cover-degree-prefix} 那一条里的 \cs{ziju} 是故意
% 不加分组的，往后管到 \cs{textbf} 那一组结束，本科封面那一行“本科”“开题报告”
% 都按这个字距排。
%    \begin{macrocode}
\ExplSyntaxOff
\hit@presetup{ term = {
  cover-university-font    = {},
  cover-document-type-font = {\erhao},
  cover-report-type-line   = {},
  cover-degree-prefix      = {},
} }
\hit@presetup[campus=weihai,degree-level=master ; campus=shenzhen,degree-level=master]{ term = {
  cover-university-font = {\ziju{1/6}},
} }
\hit@presetup[campus=shenzhen,degree-level=master,stage=proposal]{ term = {
  cover-document-type-font = {\xiaoyi\ziju{1/8}},
} }
\hit@presetup[campus=harbin,degree-level=master,stage=proposal ; campus=harbin,degree-level=master,stage=interim]{ term = {
  cover-report-type-line = {\songti\sanhao\textbf{（<report-type-zh>）}},
} }
\hit@presetup[campus=harbin ; campus=weihai]{ term = {
  cover-degree-prefix = {\ziju{1/12}\hit@term@cover@document@type@level},
} }
\ExplSyntaxOn
\ExplSyntaxOff
\hit@presetup{ term = {
  mainmatter-eqdenote-label-zh = {式中},
} }
\hit@presetup{ term = {
  mainmatter-eqdenote-label-en = {where},
} }
\hit@presetup{ term = {
  mainmatter-eqdenote-dash = {——},
} }
\hit@presetup{ term = {
  declarations-author-signature = {作者签名：},
} }
\hit@presetup{ term = {
  declarations-supervisor-signature = {导师签名：},
} }
\hit@presetup{ term = {
  declarations-date = {日期：},
} }
\hit@presetup{ term = {
  declarations-authorization = {学位论文使用权限},
} }
\hit@presetup{ term = {
  cover-school-bottom-mark-harbin-master = { 
   \begin{center} 
     \songti\sanhao\textbf{研究生院制} 
   \end{center} 
},
  cover-school-bottom-mark-harbin-doctor = { 
   \begin{center} 
     \songti\sanhao\textbf{研究生院制} 
   \end{center} 
},
} }
\hit@presetup{ term = {
  cover-school-bottom-mark-bachelor = { 
   \begin{center} 
     \hit@cover@bottommark 
   \end{center} 
},
} }
\hit@presetup{ term = {
  cover-school-bottom-mark-shenzhen-master = { 
   \begin{center} 
     \songti\sanhao\textbf{深圳校区教务部\nobreakspace{}制} 
   \end{center} 
},
} }
\hit@presetup{ term = {
  cover-note-interim-shenzhen-doctor = { 
\thispagestyle{empty} 
{\begin{center} 
\songti\sanhao\textbf{\hit@spread@to{6em}{说明}} 
\end{center} 
} 
\vspace{1cm}\songti\hit@fontsize@x{sihao}{2} 
\begin{enumerate}[itemsep=-4bp] 
\item 博士研究生学位论文中期报告一般在入学后的第五学期末完成。 
\item 此报告由学生本人填写，请导师签字后提交一份给检查小组。答辩结束后由检查组长签署意见。 
\item 报告经中期报告检查组长签字同意后，由学科部保存，以备论文答辩时参考。研究生院学生培养处将对研究生的学位论文中期报告进行抽查。 
\item 报告统一用A4纸打印。 
\end{enumerate} 
},
} }
\hit@presetup{ term = {
  declarations-date-blank = {\hspace{2.5em}年\hspace{1.5em}月\hspace{1.5em}日},
} }
\hit@presetup{ term = {
  base-hi = {嗨！thesis},
} }
\ExplSyntaxOn
%    \end{macrocode}
%
% \changes{v3.0.22}{2022/05/16}{深圳校区中期模板封面范例中没有“哈工大（深圳）制”字样。\issue[126]}
%    \begin{macrocode}
% \newcommand{\hit@shenzhen@school@bottom@mark}{
%    \begin{center}
%      \songti\sanhao\textbf{哈工大（深圳）制}
%    \end{center}
%    \begin{center}
%      \songti\sanhao\textbf{二〇一二年三月}
%    \end{center}
% }
%    \end{macrocode}
%
% 此处删除了时间
% \changes{v3.0.06}{2020/09/13}{此处删除了时间}
%
% \changes{v3.1c}{2023/04/16}{根据 \issue[151] 将本部本科开题下方的标识更换为 \texttt{.eps} 文件}
%
% \changes{v3.1c}{2023/04/16}{深圳校区硕士开题模板含有“深圳校区教务处 制”。\issue[141]}
%
%    \begin{macrocode}
%    \end{macrocode}
%
%    \begin{macrocode}
\hit@presetup@scope{}
%    \end{macrocode}
%
% \emph{勾选的那一对标签两处在用}：本科内封头一行与深圳本科报告首页那一行，
% 所以不挂档位。
%    \begin{macrocode}
\hit@presetup{ term = {
  titlepage-mark-dissertation = {毕业论文},
  titlepage-mark-practice     = {毕业设计},
} }
%    \end{macrocode}
%
%    \begin{macrocode}
%</hithesis-cfg>
%    \end{macrocode}
%
% \Finale
